\section{完成难度证明}\label{sec:completion}

参数顺序为
\[
 (t,\eps,\mu)\longrightarrow(s,m,p_*,\theta)
 \longrightarrow b\longrightarrow n\longrightarrow T\longrightarrow a.
\]
抽取常数在字母表维数 $n$ 已知之前就固定所需的博弈可靠性 $b$。
重复次数随后选定以隐藏被植入坐标，博弈的完备性误差 $a$ 最后选取。

\begin{proof}[\Cref{thm:gap} 的证明]
按定理固定 $t$ 与 $\eps$。取 $\delta=\eps$。由于
$\eps<(t-B(t))/4<(1-B(t))/4$，这满足 $0<\delta<1-B(t)$。选有理数
$0<\mu<1/2$ 使 $t\mu\le\eps/4$。
以 $t,\delta,\mu$ 应用 \Cref{prop:extraction}，固定
$s,m,p_*,\theta$，从而固定 $k,d,r$。选有理数
\[
 0<b<2^{-2s}\frac{p_*\theta}{2}.
\]
由 \Cref{prop:affine-games}，这固定字母表维数 $n$，使博弈具有任意
小的完备性误差。选整数 $T$ 满足
\[
 T\ge\frac{2^{n(2k-2s)}}{(p_*\theta)^2}.
\]
最后选正有理数 $a<1-b$ 足够小使 $2Ta\le\eps/4$，并从 NP 难的博弈
间隙 $\Val(\mathcal G)\ge1-a$ 对 $\Val(\mathcal G)\le b$ 出发。

在 YES 情形，\Cref{lem:completeness} 给出割值至少
$(1+t)/2-\eps/2$，强于所需。在 NO 情形，假设存在值至少
$(1+B(t))/2+\eps/2$ 的割。由 \Cref{lem:cut-stability}，其奇平均的
期望切片稳定性至少为 $B(t)+\eps$。
\Cref{prop:extraction,prop:decoding} 随后将蕴含
$\Val(\mathcal G)\ge2^{-2s}p_*\theta/2>b$，矛盾。故每个割的值都
小于 $(1+B(t))/2+\eps/2$。

为完备起见：该测试是确定性的多项式规模图构造。顶点集有
\[
 |X|^T\,2^{|\mathcal B|}
 =|X|^T\,2^{2^{(n+1)T}}
\]
个顶点。除 $|X|$ 外的所有因子对固定的 $t,\eps$ 都是常数。端点与边
元组只有多项式多个，且所有词生成样本空间规模为常数。枚举每个测试
结果，把其概率加到相应边权上。初始博弈权重、$t$ 与 $r$ 都是有理
数，故这些权重具有多项式位复杂度。条件边概率同样具有多项式位复杂
度。所有字母表与词都可以在相对输入规模为常数的时间内枚举。因此所
得图可在确定性多项式时间内算出。它具有有理边分布，并满足所述的、
带 $\eps$ 损失的较弱界。
\end{proof}

\begin{proof}[\Cref{thm:main} 的证明]
把 $\rho=-t$ 代入 \eqref{eq:gw-constant} 给出
\begin{equation}\label{eq:gw-ratio}
 \GW=\min_{-1<t\le1}\frac{1+B(t)}{1+t},
\end{equation}
其中同一公式对负 $t$ 定义 $B(t)=2\arcsin(t)/\pi$。
当 $-1<t\le0$ 时比值至少为 1，而当 $0<t<1$ 时小于 1。在 $t=1$ 处
等于 1，当 $t\downarrow-1$ 时发散。故其最小值在 $(0,1)$ 内取得。

固定 $\GW<\alpha\le1$。由连续性与有理数的稠密性，选有理数
$t\in(0,1)$ 使
\[
 \alpha\frac{1+t}{2}>\frac{1+B(t)}{2}.
\]
选正有理数 $\eps$ 与 $\epsilon'$ 足够小，使 $\eps<(t-B(t))/4$ 且
\begin{equation}\label{eq:strict-final-gap}
 \alpha\left(\frac{1+t}{2}-\eps-\epsilon'\right)
 >\frac{1+B(t)}{2}+\eps+\epsilon'.
\end{equation}
应用 \Cref{thm:gap}。删除自环并保留原权重尺度，使剩余权重之和至多
为 1。以误差 $\epsilon'$ 应用 \Cref{lem:unweighted}，产生简单无权
图：在公共尺度 $Q$ 上，其最大割与原最大割权重相差至多 $\epsilon'$。
在该图上，$\alpha$-近似算法在 YES 情形返回缩放值至少为
\eqref{eq:strict-final-gap} 左端的割。在 NO 情形，没有割的缩放值
超过其右端。取严格介于二者之间的固定有理阈值即可区分 NP 难的承诺
情形。定理得证。
\end{proof}
